% !TeX encoding = UTF-8

\chapter{绪论}

\section{研究背景与意义}

近年来，人工智能与大数据技术在医疗健康、金融风控与智慧城市等领域取得了突破性进展。然而，随着《中华人民共和国数据安全法》《中华人民共和国个人信息保护法》等法律法规相继实施，数据孤岛与隐私合规之间的矛盾日益凸显。传统集中式机器学习范式需要将各参与方的原始数据汇聚至中心服务器统一训练，不仅带来高昂的传输与存储开销，还存在严重的隐私泄露隐患。

为应对上述挑战，联邦学习（Federated Learning，FL）作为一种新兴的分布式协同训练范式应运而生\cite{mcmahan2017fedavg}。在联邦学习框架中，原始训练样本始终保留在本地客户端，参与方仅周期性地与中心服务器交互模型参数或梯度更新。由文献\citeinline{kairouz2021advances}可知，联邦学习在数据分布高度异构（Non-IID）与恶意推理攻击并存的开放环境下，依然面临收敛速度慢、通信代价高与隐私效用权衡困难等核心瓶颈。

\section{国内外研究现状}

\subsection{联邦聚合优化研究进展}

经典的联邦平均算法（FedAvg）通过在客户端执行多步本地随机梯度下降并按样本量加权平均，有效降低了通信频率\cite{mcmahan2017fedavg}。但在非独立同分布数据场景下，本地目标函数与全局目标函数之间的偏差会诱发显著的客户端漂移现象。在工程优化与文献著录实践中，严谨的格式规范与统计建模方法同样至关重要\cite{chenjianjun2001,gbt7714_2015}。

\subsubsection{非独立同分布数据缓解机制}

面向统计异构问题，现有研究主要沿两条技术路线展开：其一是在本地优化目标中引入正则化约束或控制变量，以校正本地梯度方向\footnote{此处演示规范 1.9 节所规定的页下“脚注”注释方式，适用于对正文术语或未公开发表资料的补充说明。}；其二是在服务器端设计基于梯度相似度或历史动量的自适应聚合策略。

\subsubsection{差分隐私与通信压缩协同设计}

已有研究表明，单纯传输梯度信息仍可能遭受梯度反演攻击并还原原始样本\cite{mahuan2011}。结合差分隐私噪声注入与梯度稀疏化能够兼顾安全性与带宽约束，相关专利与标准规范也提供了重要工程借鉴\cite{jiangxizhou1989,dublincore2012}。

\section{本文主要研究内容与结构安排}

本文严格遵循《海南大学研究生学位论文撰写规范（修订）》（海大党政办〔2023〕47 号）的结构要求，全文章节安排如下：第 1 章介绍研究背景与国内外现状；第 2 章阐述自适应聚合框架、数学公式推导与三线表及跨页续表排版示例；第 3 章展示实验配置、曲线图与对比结果分析；第 4 章总结全文主要结论并展望未来研究方向。
